\documentclass[12pt,titlepage]{article}

\usepackage{geometry}
\geometry{
    a4paper,
    total={210mm,297mm},
    left=20mm,
    right=20mm,
    top=20mm,
    bottom=20mm,
}

	
\usepackage[]{subcaption}

\usepackage{pdflscape}

% Polski
\usepackage[]{polski} 
\usepackage[polish]{babel}

%do tabel
\usepackage{multirow}

% Pierwszy akapit - wcięty
\usepackage[]{indentfirst}

% Matematyka
\usepackage[]{amsfonts}

\usepackage[]{amsmath}

% Formatowanie
\usepackage{ragged2e}

% Tytuły sekcji
\usepackage{titlesec}
%\titleformat{\section}[block]{\Large\bfseries}{}{1em}{}

% <=
\usepackage{amssymb}

% eps
\usepackage{graphicx}
% \usepackage{subfigure}

% Tabele
\usepackage{array}

\usepackage[style=czech]{csquotes}

\renewcommand*{\thesubsubsection}{}

\usepackage{hyperref}
\hypersetup{
    colorlinks,
    citecolor=black,
    filecolor=black,
    linkcolor=black,
    urlcolor=black
}

\usepackage[numbered]{matlab-prettifier}
\lstset{
    literate={ą}{{\k{a}}}1
    {Ą}{{\k{A}}}1
    {ę}{{\k{e}}}1
    {Ę}{{\k{E}}}1
    {ó}{{\'o}}1
    {Ó}{{\'O}}1
    {ś}{{\'s}}1
    {Ś}{{\'S}}1
    {ł}{{\l{}}}1
    {Ł}{{\L{}}}1
    {ż}{{\.z}}1
    {Ż}{{\.Z}}1
    {ź}{{\'z}}1
    {Ź}{{\'Z}}1
    {ć}{{\'c}}1
    {Ć}{{\'C}}1
    {ń}{{\'n}}1
    {Ń}{{\'N}}1
}

\title{
\includegraphics[scale=0.75]{img/politechnika_sl_logo_bw_poziom_pl.eps}\\
\textbf{Wydział Automatyki, Elektroniki\\
i Informatyki}\\
\vspace*{1cm}
Systemy Interaktywne i Multimedialne \\ Projekt \\ Detekcja emocji w głosie

\vspace*{5cm}
}
\author{
Natalia Stręk
Jakub Kula,
Paweł Wójtowicz,
} 
\date{Gliwice 2023}

\begin{document}

\maketitle


\newpage
\section{Analiza wyników i wnioski}
\subsection{Przedstawić najważniejsze scenariusze użycia stworzonego oprogramowania(zazwyczaj od 3 do 6 scenariuszy). Scenariusze zilustrować zrzutami ekranowymi, do których referencje mają znaleźć się w tekście.}
Wsparcie w obsłudze klienta:
Program może być zintegrowany z systemami call center, gdzie automatycznie analizuje emocje klientów podczas rozmów. Na podstawie wykrytych emocji, takich jak frustracja czy zadowolenie, system może przekazywać zgłoszenia do odpowiednich działów lub oferować automatyczne rozwiązania, poprawiając jakość obsługi klienta\\

Zarządzanie zasobami ludzkimi:
W trakcie rekrutacji i rozmów kwalifikacyjnych program może analizować emocje kandydatów, co pomoże rekruterom w ocenie autentyczności i stresu kandydatów. Może także być używany do monitorowania stanu emocjonalnego pracowników podczas sesji feedbackowych, co pozwoli na szybszą reakcję na potencjalne problemy.\\

Interaktywne systemy rozrywki:
Program może być używany w grach komputerowych i aplikacjach VR/AR, gdzie analizuje emocje graczy, aby dostosować poziom trudności gry, fabułę lub interakcje z postaciami w czasie rzeczywistym. Dzięki temu doświadczenie z gry staje się bardziej immersyjne i spersonalizowane.\\

\subsection{Do sprawozdania można dołączyć wideo prezentujące użycie scenariuszy (w osobnym pliku sprawozdania). Także w tym wypadku zespół, może być nadal zobowiązany do prezentacji wyników projektu na zajęciach.}

\subsection{Czy projekt zrealizowany został zgodnie z założeniami? Jeśli nie,podać przyczyny odstępstw od pierwotnych założeń. Należy odnieść się do zaakceptowanego dokumentu Informacje o realizowanym projekcie (ok 100 słów)}

\subsection{Czy projekt zrealizowany został zgodnie z harmonogramem? Jeśli nie, podać przyczyny opóźnień w odniesieniu do etapów pracy oraz opisać działania naprawcze podjęte aby osiągnąć wyznaczone cele projektu.}
Projekt jest realizowany zgodnie z harmonogramem

\subsection{Czy wystąpiły jakieś nadzwyczajne wydarzenia w trakcie realizacji projektu? Czy wprowadzono jakieś znaczące zmiany do zrealizowanych działań lub w składzie zespołu? Jeśli tak, opisać jakie oraz podać ich wpływ na realizację projektu.}

\subsection{Doświadczenia uzyskane podczas realizacji projektu. Opisać jaką nową wiedzę uzyskał uczestnik projektu w czasie jego realizacji np. nowe metody programowania, nowe biblioteki, nowe narzędzie itp. (ok 200 słów)}
Podczas realizacji projektu rozpoznawania emocji w głosie nauczyło się korzystania z biblioteki Librosa do szybkiej ekstrakcji cech dźwiękowych w dziedzinie częstotliwościowej oraz czasowej, co umożliwiło efektywne przetwarzanie danych audio. Opanowano również użycie biblioteki Optuna do optymalizacji hiperparametrów, co pozwoliło na zwiększenie dokładności modelu z 77\% do 80\% na zbiorze testowym. Dodatkowo, zdobyto umiejętności tworzenia graficznych interfejsów użytkownika przy pomocy biblioteki Tkinter, co umożliwiło stworzenie intuicyjnej aplikacji prezentującej wyniki analizy głosu.

\subsection{Proponowane ulepszenia projektu. Opisać możliwe zmiany w systemie, użytej technologii, narzędziach, które mogłyby prowadzić do ulepszenia projektu.}
Ulepszenie sieci neuronowej może być trudnym zadaniem, dlatego proponowane zmiany powinny skupić się na wykorzystaniu dodatkowej sieci do rozpoznawania emocji na podstawie mimiki twarzy, ponieważ mimika jest istotnym elementem komunikacji niewerbalnej. Można zastosować bibliotekę OpenCV lub FaceNet do ekstrakcji cech i klasyfikacji emocji. Integracja wyników z analizy dźwięku i obrazu w jednym systemie oraz stworzenie warstwy fuzji pozwoli na bardziej precyzyjne rozpoznawanie emocji.
\subsection{Dodać skompresowany kod źródłowy (w osobnym pliku sprawozdania)}
\end{document}